\section{Preliminaries}
\subsection{Digital signatures}

\begin{definition}[Digital signature scheme]\label{def:sig}
    A digital signature scheme contains a key generation algorithm, which outputs a secret key \sk and a public key \pk,
    and the following two algorithms:
    \begin{itemize}[leftmargin = *]
        \item
            $\opSign{\sk, m} \rightarrow \sig$:
            The signing algorithm takes as input a private key
            \sk and a message $m \in \zo^*$ and returns a signature $\sig$.
        \item 
            $\opVerify{\pk, m, \sig} \rightarrow b \in \zo$:
            The verification algorithm takes as input a public
            key $\pk$, a message $m$, and a signature $\sig$, and outputs a bit $b\in \zo$.
    \end{itemize}
\end{definition}

We say that \sig is a \emph{valid} signature on message $m$ with respect to some public key \pk if $\opVerify{pk, m, \sig} = 1$.

A digital signature scheme satisfies the \emph{correctness} and \emph{unforgeability} properties.
For all possible keypairs $(\sk, \pk)$ returned by the key-generation algorithm,
\emph{correctness} demands that $\opVerify{\pk, m, \opSign{\sk, m}} = 1$, \ewnp.
\emph{Unforgeability} demands that an adversary, who is given access to an oracle that signs messages of the adversary's choice using \sk,
cannot, \ewnp, produce a tuple $(m, \sig)$, such that \sig is a valid signature on $m$ with respect to \pk, without calling the oracle on message $m$.
